\section{Methods}
\label{sec:methods}

This review follows the PRISMA 2020 reporting standard \parencite{Page2021_prisma} and the methodological guidance of the Cochrane Handbook \parencite{Higgins2023_cochranehandbook}, adapted throughout for a synthesis of proportions rather than a synthesis of comparative treatment effects, for the reasons given in the Introduction. Every design choice below---eligibility, search, extraction, risk-of-bias assessment, and the pooling model itself---was fixed before formal data extraction began, and is reported here exactly as specified, including the places where the evidence base did not cooperate with the plan and the places where this review's own execution fell short of its own plan.

\subsection{Protocol and Registration}
\label{sec:methods-protocol}

This review was not registered in PROSPERO before literature screening or data extraction began. The design pivoted once during discovery, from a comparative-effect meta-analysis (phage-antibiotic combination versus antibiotic monotherapy) to the stratified meta-analysis of proportions reported here, after two independent search rounds found zero to one qualifying dual-arm comparative studies for MDR/XDR \textit{P.\ aeruginosa}. That pivot needed to happen before a protocol could be written that the team would not have to immediately amend. That scoping work included four rounds of literature searching (Section~\ref{sec:methods-search}) and one informal, single-reviewer screening pass of a 201-document local PDF corpus, conducted against this review's PICO criteria but without the dual-independent review PROSPERO registration presumes. PROSPERO's own policy does not prohibit preliminary scoping before registration; it prohibits registering a protocol once data extraction and synthesis are already substantively underway. By that standard, registration remained available through the design-pivot and literature-search stages of this project, and did not occur. Data extraction and the statistical synthesis reported in Results proceeded before a protocol was ever submitted -- a genuine deviation from prospective registration -- and that extraction itself was conducted by a single reviewer rather than the dual-independent process a registered protocol would specify (Sections~\ref{sec:methods-selection} and \ref{sec:methods-extraction}). The full scoping timeline is as follows: two PubMed/ClinicalTrials.gov search rounds and a targeted third and fourth round (Section~\ref{sec:methods-search}), one ad hoc single-reviewer PDF screen, and the resulting design pivot, all preceded any protocol filing, and no protocol was filed before extraction and pooling were carried out. We recommend retrospective PROSPERO registration disclosing this exact sequence before this review is submitted for publication.

\subsection{Eligibility Criteria}
\label{sec:methods-eligibility}

Eligibility follows a PICO framework, broadened partway through discovery once the design pivoted from a comparative-effect estimand to a proportion estimand (Section~\ref{sec:methods-protocol}): under the original comparative-only design, a study needed both a phage-containing arm and a within-study antibiotic-monotherapy arm to qualify, a requirement that essentially no study in this literature satisfies. Under the current design, single-arm and comparative studies are equally eligible, because every outcome is computed at the study-arm level rather than as a paired comparison.

\begin{itemize}
\item \textbf{Population.} Patients with an infection attributed to multidrug-resistant (MDR), extensively drug-resistant (XDR), or pandrug-resistant (PDR) \textit{Pseudomonas aeruginosa}, classified per the \textcite{Magiorakos2012_mdrxdr} criteria where the primary study reports sufficient susceptibility data, and per the study's own MDR/XDR/PDR label otherwise (Section~\ref{sec:methods-extraction}). Two boundary cases must be distinguished, and this review treats them differently. Where a study reports susceptibility data that \emph{positively fail} the Magiorakos threshold, the arm does not meet the population criterion and is excluded, whatever the study's own descriptive language; one arm is excluded on exactly this ground (Section~\ref{sec:results-selection}). Where a study reports \emph{no} resistance information, the arm is retained and coded \emph{not classifiable} rather than excluded, because absence of a reported susceptibility panel is not evidence of susceptibility, and excluding on missing data would discard entire trials for a reporting omission. We state plainly that this second rule makes the analysed population broader than the criterion as written: the not-classifiable arms are predominantly trial populations enrolled on the basis of \textit{P.\ aeruginosa} infection or chronic colonization without any resistance requirement, so pooled ``overall'' estimates should be read as describing phage therapy in \textit{P.\ aeruginosa} infection generally, not exclusively in the resistant population named in this review's title. Resistance-stratified estimates carry no such ambiguity. A sensitivity analysis quantifies the consequence (Section~\ref{sec:results-robustness}).
\item \textbf{Intervention.} Bacteriophage therapy, administered as monotherapy or adjunctively with antibiotics, by any route.
\item \textbf{Comparator.} None required. Where a genuine within-study comparator exists (an antibiotic-only or placebo arm), it is retained and reported, but its presence is not an eligibility condition.
\item \textbf{Outcomes.} Clinical success and any adverse event (primary); microbiological eradication, mortality, length of hospitalization, and resistance emergence during treatment (secondary), each defined per the primary study's own reporting given the expected heterogeneity in how ``success'' is defined across this literature. Because each study defines clinical success on its own terms \emph{and} across non-comparable infection syndromes --- biliary, cystic-fibrosis lung, prosthetic-joint, vascular-graft, urinary-tract, bacteremia, and osteomyelitis among those represented in this corpus --- the pooled clinical-success proportion aggregates outcomes that are not the same clinical construct, and should be read as a rough descriptive summary rather than a harmonized effect estimate. A syndrome-harmonized success definition was not achievable given the heterogeneous reporting of the primary literature, and this constraint is carried forward as a limitation (Section~\ref{sec:discussion-limitations}).
\item \textbf{Study designs.} Randomized controlled trials, prospective and retrospective cohorts, case series, and case reports, including compassionate-use reports.
\end{itemize}

A study-arm was excluded if its causative pathogen could not be separated from a mixed-pathogen cohort (coded \texttt{mixed-not-separable} in the extraction schema; Section~\ref{sec:methods-extraction}), regardless of how well the rest of the study otherwise matched. The primary electronic search (Section~\ref{sec:methods-search}) was restricted to records published from 2016 onward across all three core databases. This 2016 cutoff was chosen as a pre-specified window that comfortably predates the modern clinical phage-therapy case-report era for \textit{P.\ aeruginosa} --- which begins in earnest around 2018--2019 (the earliest included study is Duplessis et al., 2018) --- while still excluding the sparse, mostly single-institution reports from before the current compassionate-use framework. A targeted screen of the 2016--2018 window (Section~\ref{sec:methods-search}) confirmed that the pre-2019 literature for this pathogen is almost entirely preclinical (phage isolation, characterization, biofilm and animal-model work); the only extractable human clinical studies in that window are Duplessis et al.\ (2018) and the PhagoBurn trial \parencite{Jault2019_phagoburn}, both of which are included. A small number of still-older studies were also identified through citation-chasing of other reviews' reference lists (Section~\ref{sec:methods-search}) and screened against the same eligibility criteria irrespective of publication date; one such study (Rhoads et al., 2009) failed eligibility on grounds unrelated to date and is retained only as background context. The discovery-phase search separately confirmed that, contrary to the project's initial premise, at least eight further systematic reviews, meta-analyses, or scoping reviews of human phage therapy have appeared since \textcite{ElHaddad2019_eskape}, including a broader synthesis in a comparably high-impact journal \parencite{Uyttebroek2022_lancetid} and a 130-study individual-patient-data meta-analysis published only months before this review's search began \parencite{Liu2025_ijaa}.

\subsection{Information Sources and Search Strategy}
\label{sec:methods-search}

The primary electronic search was run directly against PubMed/MEDLINE using the NCBI E-Utilities API (\texttt{esearch}, \texttt{esummary}, \texttt{efetch}), and against ClinicalTrials.gov using its structured API (v2), across four documented rounds between 2026-07-18 and 2026-07-20. A representative query, run to test the central negative claim of this review (that essentially no dual-arm comparative study of phage-antibiotic combination versus antibiotic monotherapy exists for MDR/XDR \textit{P.\ aeruginosa}), took the form:

\begin{quote}
\footnotesize\ttfamily
("Pseudomonas aeruginosa"[tiab] OR "P. aeruginosa"[tiab]) AND ("phage"[tiab] OR "bacteriophage"[tiab]) AND ("multidrug-resistant"[tiab] OR "MDR"[tiab] OR "extensively drug-resistant"[tiab] OR "XDR"[tiab]) AND ("matched cohort"[tiab] OR "propensity score"[tiab] OR "registry"[tiab] OR "retrospective comparative"[tiab] OR "case-control"[tiab])
\end{quote}

\noindent This query, built directly from the terms a referee would be expected to demand, returned one hit, and that hit was a conference-abstract index record with no substantive connection to phage therapy: an effective true hit count of zero. Across all four rounds, PubMed queries of this kind returned 54 raw hits screened individually for the comparative-design question and 48 raw hits for the systematic-review/meta-analysis recheck; ClinicalTrials.gov structured queries returned 10 additional trial records screened individually. Round 3 added five previously uncatalogued comparative or randomized trials (Krakhotkin et al., 2025; Leitner et al., 2021; Karn et al., 2024; Rhoads et al., 2009; Stanley et al., 2025), found by citation-chasing the risk-of-bias tables in the largest existing phage-therapy meta-analysis's supplementary dataset \parencite{Liu2025_ijaa} instead of a new keyword search. This is standard, low-risk PRISMA practice: only citation leads were imported at this stage (Section~\ref{sec:methods-extraction} details the boundary between this citation-chasing use and any reuse of Liu et al.'s outcome data). Round 4 conducted a targeted, independent PubMed search for antibiotic-only cohort studies in MDR/XDR \textit{P.\ aeruginosa}, unrelated to any phage-therapy keyword, specifically to check whether an external comparator population could substitute for the missing within-study monotherapy arm; that search returned 44 hits and confirmed that antibiotic-only outcome cohorts are abundant in the general infectious-disease literature, but none contains a phage arm and none is therefore usable as anything other than background context (Section~\ref{sec:methods-selection}).

We also queried WHO ICTRP and PROSPERO directly. Both attempts failed for the same reason: WHO ICTRP's search portal is an ASP.NET postback-driven form, and PROSPERO's search interface returns HTTP 403, and neither can be driven by a static, unauthenticated query in the environment available to this review team. A web-search-mediated check of both registries found no competing protocol matching this review's exact population, intervention, and stratification design, but this constitutes a partial, not exhaustive, verification. A researcher with direct, authenticated access to both portals should re-run this check before this review's registration (Section~\ref{sec:methods-protocol}) is finalized.

The search rested on three databases --- PubMed/MEDLINE, ClinicalTrials.gov, and Scopus, all searched natively --- supplemented by citation-chasing, a standard and defensible base for a systematic review. To place PubMed/MEDLINE on the same topical footing as Scopus rather than leaving it represented only by the narrow comparative-design rounds above, a broad topical PubMed/MEDLINE search was run directly against the NCBI E-Utilities API on 2026-07-23, mirroring the Scopus scope: title/abstract terms for \textit{Pseudomonas aeruginosa}, for phage or bacteriophage therapy, and for multidrug, extensive, pan-, or general drug resistance, restricted to the pre-specified 2016--2026 window:

\begin{quote}
\footnotesize\ttfamily
("Pseudomonas aeruginosa"[tiab] OR "P. aeruginosa"[tiab]) AND ("phage"[tiab] OR "bacteriophage"[tiab] OR "phage therapy"[tiab] OR "bacteriophage therapy"[tiab]) AND ("multidrug-resistant"[tiab] OR "MDR"[tiab] OR "extensively drug-resistant"[tiab] OR "XDR"[tiab] OR "pandrug-resistant"[tiab] OR "PDR"[tiab] OR "drug-resistant"[tiab] OR "drug resistance"[tiab])
\end{quote}

\noindent This topical PubMed search returned 481 records over 2016--2026. Screened by document type and title against the same eligibility criteria, every eligible human clinical study it surfaced was already among those identified through the native Scopus search --- consistent with Scopus indexing a superset of this topic's clinical literature --- so the topical PubMed search added no unique-new eligible study, and is documented here as a topical core-database search confirming the Scopus-identified corpus rather than extending it. A dedicated screen of the 2016--2018 sub-window (37 of the 481 records) confirmed it contains no extractable human clinical study beyond Duplessis et al.\ (2018) and PhagoBurn \parencite{Jault2019_phagoburn}, both already included: the remaining pre-2019 records are preclinical (phage isolation, characterization, biofilm, and animal-model studies). The Scopus search was run directly against the native Scopus interface on the same date, exporting the full result set as an 860-record RIS file for screening. The query combined \texttt{TITLE-ABS-KEY} terms for \textit{Pseudomonas aeruginosa}, for phage or bacteriophage therapy, and for multidrug, extensive, pan-, or general drug resistance. The RIS export executed used \texttt{PUBYEAR > 2018}; because the topical PubMed search above establishes that the 2016--2018 window holds no additional eligible clinical study for this pathogen, the corpus is unaffected by this bound, though a Scopus re-export from 2016 is recommended before final submission for a fully uniform per-database audit trail:

\begin{quote}
\footnotesize\ttfamily
TITLE-ABS-KEY("Pseudomonas aeruginosa" AND (phage OR bacteriophage) AND (MDR OR XDR OR PDR OR "multidrug-resistant" OR "drug-resistant")) AND PUBYEAR > 2018
\end{quote}

\noindent All 860 records were screened. The native Scopus search materially enlarged the corpus: it surfaced ten eligible human clinical \textit{P.\ aeruginosa} phage-therapy studies the PubMed/MEDLINE and ClinicalTrials.gov rounds had not returned --- K\"ohler et al.\ (2023), a single inhaled-phage MDR case, and Chan et al.\ (2025), a nine-patient nebulized-phage cystic-fibrosis cohort, together with eight further reports: Law et al.\ (2019), Hahn et al.\ (2023), Lev\^eque et al.\ (2023), Duplessis et al.\ (2018), Denis et al.\ (2026), Malhotra et al.\ (2026), Yang et al.\ (2025), and \textcite{Li2025_hlife_biliary}, an 88-year-old woman with an MDR biliary-tract infection treated by multi-route phage plus antibiotics (previously carried as a pending record, now resolved from two concordant secondary sources) (Section~\ref{sec:results-selection}). It also corrected one earlier provisional exclusion: Law et al.\ (2019), previously set aside as a presumed conference abstract from the same San Diego AB-PA01 program as \textcite{Aslam2019_ajt}, is in fact a full peer-reviewed report of a distinct MDR cystic-fibrosis patient, and is now included. Four records were documented as exclusions: Van Nieuwenhuyse et al.\ (2022), a single XDR case already counted within \textcite{Pirnay2024_natmicrobiol}'s 100-case Belgian-consortium cohort (the same de-duplication basis as the \textcite{Dan2023_jid} exclusion in Section~\ref{sec:methods-selection}); Zurabov et al.\ (2023), an ICU phage-monotherapy case series whose \textit{P.\ aeruginosa} outcomes cannot be separated from a multi-pathogen cohort (otherwise one of the few monotherapy sources anywhere in this literature); Hayakawa et al.\ (2025), a multi-pathogen feasibility study with no \textit{P.\ aeruginosa}-separable outcome; and Chung et al.\ (2025), a perspective article rather than a primary study. With three databases searched --- and because Scopus indexes much of EMBASE's content --- this review meets a recognized minimum for a defensible PRISMA search. EMBASE and Web of Science were nonetheless not searched, and we state plainly why: this review team has no institutional subscription to either, so those searches could not be executed rather than merely not having been completed yet. So that any reader with subscription access can run them and audit what, if anything, they add, we report the adapted queries in full here rather than leaving them undocumented. For Ovid EMBASE:

\begin{quote}
\footnotesize\ttfamily
(pseudomonas aeruginosa or P. aeruginosa).ti,ab. AND (phage or bacteriophage or phage therapy or bacteriophage therapy).ti,ab. AND (multidrug-resistant or MDR or extensively drug-resistant or XDR or pandrug-resistant or PDR or drug-resistant or drug resistance).ti,ab. AND (2016:2026).yr.
\end{quote}

\noindent and for Web of Science Core Collection:

\begin{quote}
\footnotesize\ttfamily
TS=(("Pseudomonas aeruginosa" OR "P. aeruginosa") AND (phage OR bacteriophage OR "phage therapy" OR "bacteriophage therapy") AND ("multidrug-resistant" OR MDR OR "extensively drug-resistant" OR XDR OR "pandrug-resistant" OR PDR OR "drug-resistant" OR "drug resistance")) AND PY=(2016-2026)
\end{quote} We take two partial mitigations. First, Scopus and EMBASE overlap substantially in biomedical journal coverage, so the marginal yield of EMBASE over an already-completed native Scopus search is smaller than its absence might suggest. Second, in place of the unavailable Web of Science forward-citation function, we hand-screened the reference lists and citing literature of the three most relevant prior syntheses \parencite{ElHaddad2019_eskape,Uyttebroek2022_lancetid,Liu2025_ijaa} and ran a targeted screen of individually named landmark cases (Section~\ref{sec:results-selection}), which recovered two eligible studies no keyword search had returned. Neither mitigation is equivalent to the searches themselves, and we do not claim otherwise; the gap is a real and permanent constraint on this review, carried forward as a limitation (Section~\ref{sec:discussion-limitations}) rather than as outstanding work.

Finally, one supplementary, non-PRISMA source contributed candidate studies: an ad hoc, single-reviewer title/abstract screen of 201 PDF records assembled outside the formal database search (via a PubMed/Elicit export), screened against this review's PICO criteria using a four-criterion rubric (population, intervention, study design, and publication-date thresholds all met) replicated from an earlier 76-document reference screen. This pass returned 12 records meeting all four criteria and 189 exclusions, of which 3 of the 12 duplicated studies already in the formal search and 9 were genuinely new candidates, two of which (Krakhotkin et al., 2025 and an immunological sub-study later confirmed to duplicate an already-catalogued case) required subsequent correction (Section~\ref{sec:methods-selection}). Because this screen was conducted by a single reviewer, predates protocol registration, and did not use the dual-independent-reviewer process specified as a Field Convention for this review, we treat its disposition as a scoping input requiring independent re-screening or spot-checking, not as a component of the formal PRISMA screening record, and recommend it be re-run under the registered protocol before this review's flow diagram is finalized.

A second supplementary channel was added during peer review: a targeted landmark-case screen, prompted by a referee's observation that several canonical \textit{P.\ aeruginosa} phage-therapy cases might not have been captured by the keyword searches above. Each named landmark case was screened against this review's PICO criteria, and two proved eligible and previously uncaptured. The first is \textcite{Chan2018_omko1_emph}, the OMKO1 case: an MDR \textit{P.\ aeruginosa} chronic aortic Dacron-graft infection treated with local phage OMKO1 plus ceftazidime and resolved with no recurrence over 18 months off antibiotics (distinct from \textcite{Blasco2023_frontmed}'s vascular-graft failure and from the Chan et al.\ 2025 cystic-fibrosis cohort). The second is \textcite{Arya2026_pji_natcomms}, an MDR \textit{P.\ aeruginosa} periprosthetic joint infection treated with intermittent bacteriophage monotherapy --- no antibiotic --- over roughly two years, with clinical resolution but without microbiological eradication; it is only the second phage-monotherapy arm anywhere in this corpus (Section~\ref{sec:results-characteristics}). Both were added to the analytic dataset through this documented channel. Two further named cases were screened and not added: Rubalskii et al.\ (2020), a multi-pathogen phage series whose \textit{P.\ aeruginosa} outcomes cannot be separated from the mixed cohort (the same \texttt{mixed-not-separable} basis as other exclusions in this review), and Cano et al.\ (2021), a \textit{Klebsiella pneumoniae} prosthetic-joint case that falls outside this review's pathogen scope --- so the absence of both from the corpus is correct, not an omission. Like the PDF screen above, this landmark-case channel is a documented supplementary source, not a component of the formal database-search record.

\subsection{Study Selection}
\label{sec:methods-selection}

Screening and full-text assessment were conducted by a single reviewer at each stage (this review's data-extraction team), with a second reviewer's critique applied sequentially after each round rather than concurrently alongside the first screen. This is a deviation from the dual-independent-screening convention this field expects: no discovery-phase artifact for this project documents the personnel capacity for genuine dual-independent screening and extraction, and this review proceeded on a single-reviewer-plus-sequential-critique basis instead.

Two corrections made during screening and extraction illustrate why full-text verification, not abstract- or supplementary-table-level characterization, was necessary before any study entered the analytic dataset. First, \textcite{Krakhotkin2025_currurol} was initially catalogued as the closest structural match anywhere in this search to the monotherapy-versus-combination stratification this review set out to estimate: a four-arm, prospective observational comparative study of bacteriophage cocktail alone, bacteriophage plus one antibiotic, bacteriophage plus two antibiotics, and a two-drug antibiotic regimen with no bacteriophage, in women with multidrug-resistant chronic recurrent cystitis. It was also misclassified, in the supplementary risk-of-bias table of the largest existing phage-therapy meta-analysis, as a randomized controlled trial \parencite{Liu2025_ijaa}. Its own title, confirmed directly against the published abstract, identifies it as an observational comparative study. Full-text verification, obtained after this misclassification was caught, further confirmed that the study enrolled zero \textit{P.\ aeruginosa} cases: the causative organisms were \textit{Escherichia coli}, \textit{Klebsiella pneumoniae}, \textit{Proteus mirabilis}, \textit{Staphylococcus epidermidis}, and \textit{Staphylococcus aureus} only. Krakhotkin et al.\ (2025) was excluded from this review's evidence base in its entirety on this basis; it appears here only as an example of why this review's own screening caught an error present in a prior team's dataset, and should not be mistaken for a retained data source.

Second, an immunological sub-study reporting host T-cell and antibody responses to bacteriophage in a lung-transplant recipient on adjunctive phage therapy for multidrug-resistant \textit{P.\ aeruginosa} pneumonia \parencite{Dan2023_jid} was confirmed, on full-text comparison, to describe the identical clinical case already reported as Patient 1 in \textcite{Aslam2019_ajt}: the same bilateral lung-transplant recipient, the same two treatment episodes, the same intravenous-plus-nebulized phage regimen. This study was excluded as a confirmed duplicate, avoiding double-counting a single patient's outcome twice in the pooled estimates.

A third finding from Round 4's targeted search (Section~\ref{sec:methods-search}) concerns \textcite{Leitner2021_lancetid_pyophage}, the single cleanest genuine antibiotic-monotherapy comparator arm identified anywhere in this literature: a three-arm randomized trial of intravesical bacteriophage versus placebo versus open-label systemic antibiotics, in which \textit{P.\ aeruginosa} is confirmed among the enrolled uropathogens. Despite an exhaustive attempt to obtain a \textit{Pseudomonas}-specific subgroup breakdown, using direct PubMed full-text retrieval, a PMC linkout check, a ClinicalTrials.gov structured-results query, and a search for secondary sources citing the trial in enough detail to reconstruct a per-pathogen breakdown, no such breakdown could be located. This trial's outcome data are reported only as a six-pathogen aggregate with no route toward isolating the \textit{P.\ aeruginosa} subgroup from the sources available to this review team. It is retained in the extraction dataset as a documented, access-blocked exclusion -- an open question, not a settled negative -- since full-text or author-correspondence access, if obtained, could still resolve this study into the monotherapy stratum this review otherwise could not populate (Section~\ref{sec:methods-outcomes}).

\subsection{Data Extraction}
\label{sec:methods-extraction}

The unit of extraction is the study-arm, not the study: a study reporting outcomes for more than one regimen (for example, more than one administration route within the same cohort) contributes one row per arm. Extraction was conducted by a single reviewer per study, again without concurrent dual-independent extraction, for the same resourcing reason disclosed in Section~\ref{sec:methods-selection}. The extraction schema, fixed before extraction began, records for each study-arm: sample size; pathogen scope (\textit{Pseudomonas}-only, mixed-with-separable-subgroup, or mixed-not-separable, the last excluded per Section~\ref{sec:methods-eligibility}); resistance classification and its source; route of phage administration; treatment modality; the numerator and study-specific definition for clinical success; the numerator for any adverse event; microbiological eradication; mortality; length of hospitalization; resistance emergence during treatment; study design; publication year; journal tier; geographic source; and a documented primary-source citation locator (page, table, or verbatim quote) for every extracted value.

Resistance classification followed a four-step ladder applied in order: (1) if the primary study reports a complete antibiogram sufficient to apply the \textcite{Magiorakos2012_mdrxdr} criteria directly (non-susceptibility to at least one agent in at least three antimicrobial categories for MDR, to all but two or fewer categories for XDR, or to every agent tested in every category for PDR), the arm is coded \texttt{independently-verified}; (2) if the study labels the population MDR, XDR, or PDR under its own stated criteria without a full panel, the arm is coded \texttt{author-reported} and the author's label is used as given; (3) if only a partial panel is reported, an asymmetric, deliberately conservative rule applies: an MDR call is permitted from partial data, but an upgrade to XDR requires positive confirmation of susceptibility across enough categories to rule out a milder classification, not merely the absence of testing in the untested categories; (4) if no classification information is available at all, the arm is coded \texttt{not classifiable} and retained as its own explicit category, not merged into MDR or dropped. This asymmetric rule has a mechanical consequence: it shrinks the XDR stratum relative to a symmetric rule, because any partially tested isolate that might plausibly be XDR defaults to MDR or to not-classifiable absent explicit confirmation. This is a deliberate trade-off between classification conservatism and statistical power in the XDR stratum, and a reader comparing this review's XDR estimate against another synthesis's should attribute any difference to this coding choice before attributing it to a clinical difference.

Because the largest existing phage-therapy meta-analysis \parencite{Liu2025_ijaa} makes a case-level extraction dataset publicly available as a supplementary file, this review adopted an explicit, tiered reuse policy, stated up front. Tier 1, using that dataset purely to identify candidate studies for citation-chasing (Section~\ref{sec:methods-search}), was used as described above and is standard PRISMA practice, since no outcome data were imported at that stage. Tier 2 calls for independently cross-validating this review's own extracted values against Liu et al.'s corresponding entries for every study appearing in both datasets, with a documented primary-source citation required for a sample of ``independently extracted'' rows to guard against a coder simply copying Liu et al.'s values and mislabeling them as independent. This was adopted as the intended default policy but has not yet been executed as a formal cross-validation exercise on this extraction round: every row in the final analytic dataset is coded \texttt{data\_provenance = independently-extracted}, and no row in the final analytic dataset carries a \texttt{cross-validated-against-Liu2025} label. The planned side-by-side comparison against Liu et al.'s codings for overlapping studies remains outstanding. Tier 3, direct adoption of Liu et al.'s codings in place of independent extraction, was never invoked for any study in this dataset, consistent with the policy that Tier 3 requires explicit, per-instance approval and is expected to be rare.

The final analytic dataset contains 31 study-arms across 26 studies. One arm \parencite{Leitner2021_lancetid_pyophage} is excluded from all pooling, for the access-blocked reason given in Section~\ref{sec:methods-selection}, leaving 29 pooling-eligible study-arms across 25 studies. The largest single source, \citeauthor{Pirnay2024_natmicrobiol} (\citeyear{Pirnay2024_natmicrobiol}), reports 49 \textit{P.\ aeruginosa} patients across a 100-patient multinational cohort. An initial extraction pass captured only an 11-patient subset given full case-narrative treatment in the source paper's main text; this was superseded during drafting by a full re-extraction from the paper's open-access supplementary tables (PMC11153159), which report structured per-patient data for all 100 cases. Of the 49 \textit{P.\ aeruginosa}-targeted patients, 17 are classified ``usual drug resistance'' in the source data and do not meet this review's MDR/XDR/PDR eligibility criterion; they are excluded. The remaining 32 patients are extracted as three resistance-stratified arms (7 XDR, 23 MDR, 2 PDR), replacing the original 11-patient aggregate: all 32 eligible patients are now included, resolving the within-study selection-bias risk the narratively-selected 11-patient subset had raised. Cross-validation against that original subset found exact agreement for 9 of 11 patients; the other 2 (case \#30 and \#71) are now excluded as usual-drug-resistance, a correction made possible only once per-patient resistance data became available, reflecting new information rather than a flaw in the original extraction. A handful of arms have missing values for individual outcomes while the rest of the row remains intact: most notably, clinical success is not extractable as a binary proportion for two RCTs, whose primary endpoint is a continuous or time-to-event measure instead of a fixed-timepoint success count, and likewise for Chan et al.'s (2025) cystic-fibrosis cohort, whose endpoints are continuous ppFEV$_1$ and sputum-CFU measures rather than a binary success count; these arms are excluded from that one outcome's pooling and retained in the dataset for every other outcome.

\subsection{Risk-of-Bias Assessment}
\label{sec:methods-rob}

The pre-specified plan assigns ROBINS-I \parencite{Sterne2016_robinsi} to non-randomized comparative studies, Cochrane RoB2 \parencite{Sterne2019_rob2} to randomized controlled trials, and the JBI critical appraisal tool for case series \parencite{Munn2020_jbicaseseries} to single-arm case series and cohorts, with the case-report-specific approach of \textcite{Murad2018_casereports} reserved for the smallest reports (one to three patients), since several JBI case-series items presume a scale of reporting that does not apply at that size. ROBINS-I was not applied to any study-arm in the final analytic dataset: every non-randomized source in this corpus is single-arm (a case report, case series, or cohort without a within-study comparator), so none meets ROBINS-I's own scope of non-randomized \emph{comparative} studies; all such sources were rated with the JBI or Murad et al.\ tools instead.

Risk-of-bias ratings are now assigned for all 31 study-arms. The single-patient case reports and small case series were rated at Low-to-moderate risk using the Murad et al.\ framework; this group was enlarged by the native Scopus search, which added seven single-patient case reports rated under the same tool (Law et al.\ 2019, Lev\^eque et al.\ 2023, Duplessis et al.\ 2018, Denis et al.\ 2026, Malhotra et al.\ 2026, Yang et al.\ 2025, and Li et al.\ 2025) and Hahn et al.'s (2023) two-patient cystic-fibrosis series rated Low-to-moderate on the JBI checklist, alongside K\"ohler et al.'s (2023) single inhaled-phage case rated earlier under the same Murad et al.\ tool. The two landmark cases added during peer review --- \textcite{Chan2018_omko1_emph} and \textcite{Arya2026_pji_natcomms}, both single-patient reports --- were likewise rated Low-to-moderate risk under the Murad et al.\ tool. A per-domain tabulation of these single-patient judgments was not retained in a stored table, a disclosed, low-consequence gap given these are the tool's simplest, lowest-failure-mode cases (case reports of 1-3 patients). The remaining multi-patient and trial arms --- the four randomized controlled trials, \citeauthor{Onallah2023_med}'s (\citeyear{Onallah2023_med}) Israeli Phage Therapy Center case series, \citeauthor{Pirnay2024_natmicrobiol}'s three resistance-stratified arms, and Chan et al.'s (2025) two-arm cystic-fibrosis case series (a 9-patient consecutive series rated Low-to-moderate risk on the JBI checklist) --- were rated in a dedicated pass using RoB2 for the trials and the JBI case-series checklist for the cohort/case-series sources, drawing on structured PubMed abstracts, a direct ClinicalTrials.gov API v2 results-module query for the SWARM-P.a./AP-PA02 trial, and \citeauthor{Pirnay2024_natmicrobiol}'s full per-patient supplementary data already read in full for this review's resistance-stratified re-extraction (Section~\ref{sec:results-characteristics}). Full per-domain judgments and rationale are reported in \texttt{quality\_reports/risk\_of\_bias\_assessment\_phage\_therapy\_mdr\_pseudomonas.md}. In summary: \textcite{Pirnay2024_natmicrobiol}'s cohort rates Low risk on the JBI checklist (a particular strength: explicitly consecutive, complete, and fully transparent per-patient reporting that this review team verified firsthand); the Israeli case series rates Low-to-moderate risk (two Unclear items reflecting abstract-only access, not an apparent design flaw); the SWARM-P.a./AP-PA02 trial rates Low risk on RoB2 (the only RCT in this corpus whose completion and outcome-ascertainment domains this review confirmed directly from ClinicalTrials.gov's structured trial-registry results); \citeauthor{Weiner2025_natcomms}'s (\citeyear{Weiner2025_natcomms}) BX004-A trial and \citeauthor{Leitner2021_lancetid_pyophage}'s (\citeyear{Leitner2021_lancetid_pyophage}) trial both rate Some concerns (small, incompletely-reported safety population for BX004-A; a partially open-label comparator arm and moderate attrition for Leitner et al.); and the PhagoBurn trial rates High risk, driven by unmasked clinicians and early stopping for insufficient efficacy --- a material caveat for a trial that otherwise anchors this review's phage-monotherapy-adjacent evidence. This risk-of-bias picture is the input to the GRADE certainty assessment described below, which was carried out once it was complete.

The planned reuse of Liu et al.'s (2025) existing risk-of-bias assessments, which cover RoB2 ratings for 9 randomized trials, ROBINS-I ratings for 2 comparative cohorts, and NHLBI/JBI ratings for approximately 90 case series and reports in their own dataset, follows the same tiered logic described in Section~\ref{sec:methods-extraction}: at minimum a 20\% random sample of overlapping studies, or all studies falling within this review's three pre-specified strata, whichever is larger, was to be independently re-rated as a check before adopting any of Liu et al.'s existing ratings. As with the outcome-data cross-validation, this spot-check has not yet been performed on the current extraction round; no study-arm in the current dataset carries an \texttt{adopted-from-Liu2025-spot-checked} disposition. GRADE certainty of evidence \parencite{Guyatt2011_grade} was assessed for every pooled outcome-stratum cell. Because every cell is an uncontrolled single-arm proportion drawn predominantly from compassionate-use case reports and case series --- the randomized trials contribute only their phage arms, stripped of their comparators --- each body of evidence starts at \emph{Low} certainty, the GRADE default for observational evidence, and is then rated down across the five standard domains of risk of bias, inconsistency, indirectness, imprecision and publication bias \parencite{Guyatt2011_grade}. To keep the assessment reproducible rather than a set of unaudited per-cell judgments, the domains are applied by pre-specified rules to the cell's own data (\texttt{scripts/R/10\_grade.R}): risk of bias, rated down two levels where single-patient case reports contribute more than half the cell's patients \emph{and} the one High-risk-of-bias trial also contributes, and one level otherwise; inconsistency, rated down two levels where the 95\% prediction interval spans at least 80 percentage points and one level at 50 or more (the prediction interval rather than $I^2$, which is degenerate for a binomial-normal GLMM and is documented as frequently uninformative and misinterpreted in proportion syntheses; \cite{Migliavaca2022_i2prevalence}); indirectness, rated down one level for the compassionate-use salvage population common to every cell and two for clinical success, whose definition varies by study across non-comparable syndromes; imprecision, rated down two levels where the 95\% confidence interval spans at least 50 percentage points and one at 30 or more; and publication bias, rated down one level in every cell, since a compassionate-use case report describing a failure is markedly less likely to be written and published than one describing a success, and Peters' test cannot exclude this at the present corpus size.

Three of these five domains --- risk of bias, indirectness, and publication bias --- deduct at least one level in every cell by construction, so the minimum possible deduction is three levels against a starting certainty of Low. \emph{Every cell therefore rates Very low before the data-driven domains are consulted}, and we report the rating as deterministic rather than presenting a uniform floor as though it were a discriminating finding. Only inconsistency and imprecision vary across cells; the resulting domain profile, not the rating, is what distinguishes one cell from another (Section~\ref{sec:results-grade}). No cell qualified for the observational upgrade criteria (large magnitude, dose-response, or confounding that would work against the observed direction), none of which is estimable without a comparator. Certainty cannot fall below \emph{Very low}, GRADE's floor.

\subsection{Outcomes and Statistical Analysis}
\label{sec:methods-outcomes}

Four outcome families are pooled where the evidence permits: clinical success and any adverse event (primary), and microbiological eradication and mortality (secondary); length of hospitalization and resistance emergence during treatment were extracted but are reported narratively rather than pooled, given their sparse and heterogeneous reporting across studies. Each outcome is pooled separately within every cell of a pre-specified three-way stratification: resistance category (MDR, XDR, PDR, or not classifiable), route of phage administration (topical/local, intravenous, inhaled/nebulized, oral, intra-articular, or other), and treatment modality (phage monotherapy, phage-antibiotic combination, or antibiotic monotherapy). An unstratified, pooled-across-everything estimate is also reported for each primary outcome, explicitly labeled secondary and contextual, for comparability with prior non-stratified syntheses \parencite{ElHaddad2019_eskape,Uyttebroek2022_lancetid,Liu2025_ijaa}. This review's contribution is the stratified breakdown, not a single headline proportion.

The primary pooling model is a random-effects binomial-normal generalized linear mixed model (GLMM) with a logit link \parencite{Hamza2008_binomial,Nyaga2014_metaprop}, implemented in R (version 4.6.1) using the \texttt{meta} and \texttt{metafor} packages:
\[
r_i \mid p_i \sim \text{Binomial}(n_i, p_i), \qquad \text{logit}(p_i) = \mu_s + u_i, \qquad u_i \sim N(0, \tau_s^2),
\]
with the pooled proportion for stratum $s$ recovered as $\hat p_s = \operatorname{expit}(\hat\mu_s)$ and its 95\% confidence interval computed on the logit scale before back-transformation. We chose this model over the Freeman-Tukey double-arcsine transformation specified as the field's default variance-stabilizing approach for proportions near the 0/1 boundary \parencite{FreemanTukey1950,Miller1978_inversearcsine}, because this corpus spans study-arms as small as a single case report and as large as a 100-patient cohort, and \textcite{Schwarzer2019_doublearcsine} document that double-arcsine back-transformation under exactly this kind of sample-size heterogeneity can produce a pooled point estimate lying outside the range of every observed study. The GLMM avoids this failure mode by modeling the exact binomial likelihood at each study-arm's own sample size, with no transformation or back-transformation step. Double-arcsine (DerSimonian-Laird $\tau^2$) and a simple logit transformation with DerSimonian-Laird pooling \parencite{DerSimonianLaird1986} are reported alongside the primary model as sensitivity comparisons, and the Hartung-Knapp-Sidik-Jonkman adjustment \parencite{IntHout2014_hksj} is applied to every stratum's confidence interval given the small number of studies expected per cell \parencite{Barker2021_proportionalma}.

Because a binomial-normal GLMM fit by numerical quadrature can fail to converge in exactly the sparse, boundary-proportion regime this corpus produces (many single-patient case reports with zero or one event, and strata where every arm reports 0\% or 100\% of an outcome), every stratum-outcome cell's GLMM fit was checked for convergence (captured errors, known non-convergence warning strings, or a non-finite pooled estimate or standard error), and this diagnostic is reported for every cell, not only the cells that failed. Where GLMM failed to converge for a specific cell, the logit-DerSimonian-Laird sensitivity model was promoted to serve as that cell's reported primary estimate; this substitution is cell-level only, never model-wide, and is always disclosed with an explicit footnote naming the fallback and the reason, so that a fallback estimate is never visually indistinguishable from a converged GLMM estimate. Where a single study contributed more than one arm to the same stratum-outcome cell, robust variance estimation clustered by study (\texttt{clubSandwich}) was applied alongside the naive-independence estimate, since arms from the same study share a site, a measurement protocol, and often the same investigators. A stratum-outcome cell is reported as a formal pooled estimate only if it contains at least 3 independent studies and at least 20 pooled patients; cells below this pre-specified threshold are reported narratively instead: individual study-level proportions are tabulated with an explicit statement that the evidence is insufficient to support formal pooling, so an unstable estimate from one or two studies is never presented as though it were a stable finding. Between-stratum comparisons use a subgroup meta-analysis with Cochran's $Q$-between test on the transformed proportion scale, not a meta-regression on a paired treatment effect, since no paired effect size exists under this design; no formal correction for multiple comparisons (Bonferroni, Benjamini-Hochberg) is applied across the resulting battery of subgroup tests, a deliberate choice consistent with the Cochrane Handbook's treatment of pre-specified subgroup analyses as hypothesis-generating rather than confirmatory \parencite{Higgins2023_cochranehandbook}, since the stratification variables were fixed for clinical reasons before extraction, not chosen after seeing which splits looked significant.

Where a genuine within-study phage-versus-antibiotic-monotherapy comparator exists, a paired risk ratio was planned as a clearly labeled exploratory supplementary analysis, never blended into the primary stratified-proportions estimand. In the event, no such comparison is reported: Krakhotkin et al.\ (2025) is excluded entirely (Section~\ref{sec:methods-selection}), Leitner et al.\ (2021) is excluded from pooling for lack of a \textit{Pseudomonas}-specific breakdown, and the one remaining candidate (PhagoBurn) reports a time-to-event, not a binary, primary outcome. This exploratory analysis was planned but ultimately not reportable given the data this review could extract.

Robustness was assessed through thirteen pre-specified checks: transformation and model choice (GLMM versus double-arcsine versus logit-DerSimonian-Laird versus fixed-effect); restricting resistance classification to independently verified cases only; fixed- versus random-effects pooling; a study-design-stratified comparison (randomized and prospective designs versus retrospective and compassionate-use designs); a geographic-concentration check excluding the two most accessible cohorts (Belgium, Israel) to test whether the pooled estimate is an artifact of which data happened to be easiest to obtain; leave-one-out removal of each study; a collapsed two-category route taxonomy (systemic versus local) alongside the full route taxonomy; a transparency check comparing the pooled estimate using only independently extracted rows against all rows, given the Liu et al.\ reuse policy in Section~\ref{sec:methods-extraction}; Hartung-Knapp-Sidik-Jonkman versus standard Wald confidence intervals; robust-variance clustering for same-study multi-arm cells; a funnel plot and Peters' regression test for small-study effects---Peters' test, regressing the event proportion on the inverse sample size, rather than Egger's linear regression, because for a meta-analysis of proportions the standard error is a deterministic function of the proportion itself and manufactures funnel asymmetry near the 0/1 boundary, so Egger's test is invalid here and Peters' is the recommended small-study test for binary/proportion outcomes \parencite{Peters2006_publicationbias}---applied only to stratum-outcome cells with at least 10 studies per the pre-specified threshold, and explicitly not run, with that reason stated, below it; an appendix-only relaxation of the minimum-pooling threshold to at least 2 studies and 10 patients, to show whether the primary threshold is itself consequential without promoting the relaxed-threshold estimates to primary results; and application of the pre-specified monotherapy-stratum decision rule once full-text access to the candidate comparator studies was exhausted.

Six falsification checks, adapted from causal-inference negative-control logic to this descriptive design, were pre-specified to test whether the pooled proportions reflect a genuine, interpretable summary of the evidence rather than an artifact of reporting drift or selective inclusion: a meta-regression of the transformed clinical-success proportion on publication year, expecting a null or weak coefficient absent a documented mechanism; a comparison of pooled proportions across journal tiers, since a strong tier gradient would itself suggest editorial or selective-reporting distortion rather than a clinical finding; a visual small-study/sample-size check plotting each arm's proportion against $1/\sqrt{n}$, alongside the formal Peters' regression test; a cross-tabulation of route by resistance class by geographic source, to rule out an apparent route effect that is actually a single-center signature; a comparison of length-of-hospitalization against clinical-success and eradication, on the logic that if every outcome, including one with no obvious direct phage mechanism, moves favorably to a similar degree, that pattern is more consistent with a uniform reporting or selection artifact than with a phage-specific signal; and a check of adverse-event rates against publication year, to rule out apparent safety improvement that is actually a change in adverse-event ascertainment practice over time. The length-of-hospitalization comparison is reported as not applicable in the current analysis: only 1 of the 29 pooling-eligible study-arms reports a length-of-hospitalization value, too sparse to support even a qualitative directional comparison against the clinical-success or eradication estimates.
